\documentclass[a4paper, 12pt]{article}

\usepackage{utils}

\renewcommand*{\today}{02 April 2026}

\begin{document}

\hotbox{Analyse 4}{CM 14}{\today}

\subsection{Approximation ponctuelle}

\begin{theoreme}{Théorème de Dirichlet}{}
    Soit $f: [0, 2\pi[ \to \mathbb{C}$ une fonction $\mathcal{C}^1$ par morceaux, $2\pi$-périodique.
    Alors pour tout $x \in \R$ on a
    $$
    \lim_{n \to +\infty} S_{f,n}(x) = \dfrac{f(x_+) + f(x_-)}{2}
    $$
    où $f(x_+)$ et $f(x_-)$ sont respectivement les limites à droite et à gauche de $f$ en $x$.
\end{theoreme}

\begin{remarque}
    À l'intérieur d'un intervalle de continuité de $f$, la série de Fourier de $f$ converge vers $f$.
    
    Notre théorème nous sert quand on est en présence d'une discontinuité de $f$.
\end{remarque}

\begin{demonstration}
    \n
    \n
    \n
    \n
    \n
    \n
    \stump{À faire}{vite!}
    \n
    \n
    \n
    \n
    \n
\end{demonstration}

\begin{definition}
    On appelle \textbf{moyenne de Cesàro} d'ordre $n$ de la série de Fourier de $f$, le polynome trigonométrique
    $$
    \sigma_{f,n} = \frac{1}{n + 1} \sum_{k=0}^n S_{f,k}
    $$
\end{definition}

\begin{theoreme}{Théorème de convergence uniforme de Fejér}{}
    Si $f: \R \to \C$ est $2\pi$-périodique et continue sur $[0, 2\pi]$, alors
    $$
    \lim_{n \to +\infty} \norm{\sigma_{f,n} - f}_{\infty, [0, 2\pi]} = 0
    $$
    
    c'est à dire que la moyenne de Cesàro de la série de Fourier de $f$ converge uniformément vers $f$.
\end{theoreme}

\begin{remarque}
    Quand c'est le cas on peut alors utiliser l'identité de Parseval
\end{remarque}

\begin{demonstration}
    \n
    \n
    \n
    \n
    \n
    \n
    \stump{À faire}{vite!}
    \n
    \n
    \n
    \n
    \n
\end{demonstration}

\end{document}